\chapter{Common Mathematical and Simulation Framework}

\section{Indices, coordinate frames, and planar assumption}
All discrete-time quantities are indexed by $k\in\{0,1,\ldots\}$, and consecutive samples are separated by the sampling interval $\Delta t$. Continuous time is denoted by $t$ when a trajectory is specified analytically.

Two coordinate frames are used. The navigation frame $n$ is the fixed two-dimensional world frame in which global positions and velocities are evaluated. The body frame $b$ is rigidly attached to the target node and its inertial measurement unit (IMU). The yaw angle $\psi$ describes the rotation of the body frame relative to the navigation frame. A horizontal vector expressed in the body frame is mapped to the navigation frame by
\begin{equation}
R_b^n(\psi)=
\begin{bmatrix}
\cos\psi & -\sin\psi\\
\sin\psi & \cos\psi
\end{bmatrix},
\end{equation}
where $R_b^n(\psi)$ is the body-to-navigation rotation matrix.

Phase 1 deliberately uses a level planar model. Roll and pitch are assumed sufficiently small or already compensated, and the two horizontal acceleration channels are treated as gravity-compensated. This is an important modeling assumption. A real three-axis accelerometer measures \emph{specific force}, not directly horizontal translational acceleration. In a full three-dimensional strapdown inertial navigation system (INS), attitude must be estimated and gravity must be removed before acceleration is integrated; see, e.g., Groves~\cite{groves2013principles}. The planar model is therefore a controlled first step rather than the final hardware model.

\section{Navigation state}
The target state at sample $k$ is
\begin{equation}
\bm{x}_k=
\begin{bmatrix}
p_{x,k} & p_{y,k} & v_{x,k} & v_{y,k} & \psi_k
\end{bmatrix}^{\top}.
\label{eq:state}
\end{equation}
The individual state variables are:
\begin{itemize}
    \item $p_{x,k}$: global target position along the navigation-frame $x$-axis;
    \item $p_{y,k}$: global target position along the navigation-frame $y$-axis;
    \item $v_{x,k}$: target velocity component along the navigation-frame $x$-axis;
    \item $v_{y,k}$: target velocity component along the navigation-frame $y$-axis;
    \item $\psi_k$: target yaw angle, i.e. orientation of the body frame relative to the navigation frame.
\end{itemize}
For compact notation, the position and velocity vectors are
\begin{equation}
\bm{p}_k=
\begin{bmatrix}p_{x,k}\\p_{y,k}\end{bmatrix},
\qquad
\bm{v}_k=
\begin{bmatrix}v_{x,k}\\v_{y,k}\end{bmatrix}.
\end{equation}

\section{Ideal planar IMU signal}
The ideal planar IMU input at sample $k$ is written as
\begin{equation}
\bm{u}_k=
\begin{bmatrix}
a_{x,k}^{b} & a_{y,k}^{b} & \omega_{z,k}
\end{bmatrix}^{\top},
\end{equation}
where $a_{x,k}^{b}$ is the body-frame longitudinal acceleration, $a_{y,k}^{b}$ is the body-frame lateral acceleration, and $\omega_{z,k}$ is the angular rate about the vertical/body $z$-axis, i.e. the yaw rate. The ideal horizontal body-frame acceleration vector is
\begin{equation}
\bm{a}_k^b=
\begin{bmatrix}a_{x,k}^{b}\\a_{y,k}^{b}\end{bmatrix}.
\end{equation}

\section{IMU measurement and bias model}
The simulated accelerometer measurement $\bm{a}_{m,k}^{b}$ and gyroscope measurement $\omega_{m,k}$ are modeled as
\begin{align}
\bm{a}_{m,k}^{b} &= \bm{a}_{k}^{b}+\bm{b}_{a,k}+\bm{n}_{a,k},\\
\omega_{m,k} &= \omega_{z,k}+b_{g,k}+n_{g,k}.
\end{align}
Here, $\bm{b}_{a,k}$ is the additive two-dimensional accelerometer bias, $b_{g,k}$ is the additive gyroscope yaw-rate bias, $\bm{n}_{a,k}$ is accelerometer white noise, and $n_{g,k}$ is gyroscope white noise. In the Phase-1 simulation,
\begin{align}
\bm{n}_{a,k} &\sim \mathcal{N}(\bm{0},\sigma_a^2 I_2),\\
n_{g,k} &\sim \mathcal{N}(0,\sigma_{\omega}^2),
\end{align}
where $\sigma_a$ is the per-sample accelerometer-noise standard deviation, $\sigma_{\omega}$ is the per-sample gyroscope-noise standard deviation, and $I_2$ is the $2\times2$ identity matrix.

The sensor biases evolve as discrete random walks,
\begin{align}
\bm{b}_{a,k+1} &= \bm{b}_{a,k}+\bm{w}_{ba,k},\\
b_{g,k+1} &= b_{g,k}+w_{bg,k},
\end{align}
with increments
\begin{align}
\bm{w}_{ba,k} &\sim \mathcal{N}(\bm{0},\sigma_{ba}^2\Delta t\,I_2),\\
w_{bg,k} &\sim \mathcal{N}(0,\sigma_{bg}^2\Delta t).
\end{align}
Thus, $\sigma_{ba}$ and $\sigma_{bg}$ are continuous-time-style random-walk intensities, and the standard deviation of one bias increment is proportional to $\sqrt{\Delta t}$. The model is intentionally simpler than a complete IMU error model. Scale factors, axis misalignment, temperature effects, nonlinearity, vibration, and explicit three-dimensional gravity compensation are postponed until hardware characterization provides evidence that they are needed.

\section{Planar inertial mechanization}
Acceleration is measured in the body frame, whereas position and velocity are maintained in the navigation frame. Therefore orientation information from the gyroscope is required before the accelerometer signal can be integrated. In this report, ``acceleration-based dead reckoning'' consequently means IMU-based dead reckoning using accelerometer and gyroscope measurements.

For the interval from sample $k$ to $k+1$, the implementation first forms the midpoint yaw
\begin{equation}
\psi_{k+\frac12}=\psi_k+\frac12\omega_{m,k}\Delta t,
\end{equation}
where $\psi_{k+1/2}$ approximates the orientation halfway through the integration interval. The measured body acceleration is then rotated into the navigation frame,
\begin{equation}
\bm{a}_{k}^{n}=R_b^n\!\left(\psi_{k+\frac12}\right)\bm{a}_{m,k}^{b},
\end{equation}
where $\bm{a}_{k}^{n}$ denotes horizontal acceleration expressed in the navigation frame.

The discrete inertial mechanization is
\begin{align}
\bm{p}_{k+1} &= \bm{p}_k+\bm{v}_k\Delta t+\frac12\bm{a}_{k}^{n}\Delta t^2,\label{eq:mechanization_p}\\
\bm{v}_{k+1} &= \bm{v}_k+\bm{a}_{k}^{n}\Delta t,\label{eq:mechanization_v}\\
\psi_{k+1} &= \operatorname{wrap}\!\left(\psi_k+\omega_{m,k}\Delta t\right).\label{eq:mechanization_psi}
\end{align}
The function $\operatorname{wrap}(\cdot)$ maps angles to $[-\pi,\pi)$. The same equations, supplied with ideal noise-free IMU signals, are used to generate the deterministic Phase-1 reference trajectory. Consequently, ``ideal IMU dead reckoning reconstructs the ground truth exactly'' is a software consistency invariant: it checks indexing, frame conventions, and integration logic rather than representing an empirical sensor-performance result.

\section{Why inertial dead reckoning drifts}
The fundamental reason inertial dead reckoning drifts is integration of measurement error. Accelerometer noise and bias are integrated once into velocity and twice into position. As a simple illustration, consider a constant scalar acceleration bias $b_a$ acting along one fixed direction. Neglecting all other errors and assuming zero initial error, the induced position error $\delta p(t)$ grows approximately as
\begin{equation}
\delta p(t)\approx\frac12 b_a t^2.
\end{equation}
A gyroscope error is additionally important because it changes the estimated rotation $R_b^n(\psi)$ and therefore projects the measured body acceleration into the wrong global direction. In a full three-dimensional implementation, attitude error also contaminates gravity removal and can dominate the drift. Ultra-wideband measurements are introduced to provide external geometric information before these inertial errors become unacceptably large.

\section{UWB measurement model}
Let auxiliary node $A$ have known navigation-frame position
\begin{equation}
\bm{p}_{A,k}=
\begin{bmatrix}p_{A,x,k}\\p_{A,y,k}\end{bmatrix}
\end{equation}
at sample $k$. The ideal geometric target-to-auxiliary range is
\begin{equation}
d_k=\lVert\bm{p}_k-\bm{p}_{A,k}\rVert_2,
\end{equation}
where $d_k$ is a scalar distance. The measured UWB range is
\begin{equation}
z_k=d_k+\nu_k,
\qquad
\nu_k\sim\mathcal{N}(0,\sigma_{\mathrm{UWB}}^2),
\label{eq:uwb_measurement}
\end{equation}
where $z_k$ is the UWB measurement, $\nu_k$ is additive range noise, and $\sigma_{\mathrm{UWB}}$ is its standard deviation. The symbol $\nu_k$ is deliberately used for UWB noise so that it cannot be confused with the velocity vector $\bm{v}_k$. Packet loss, non-line-of-sight bias, outliers, and time-varying measurement quality are intentionally excluded from Phase 1A and are introduced in later experiments.

\section{Bootstrap particle filter}
A conventional bootstrap particle filter (PF)~\cite{arulampalam2002tutorial,ristic2004beyond} is used as the transparent baseline before reproducing the AACOPF-specific particle transition of Han et al.~\cite{han2020cooperative}. The PF approximates the posterior state density by
\begin{equation}
p(\bm{x}_k\mid z_{1:k},\bm{u}_{1:k})\approx
\sum_{i=1}^{N_p}w_k^{(i)}\,\delta\!\left(\bm{x}_k-\bm{x}_k^{(i)}\right),
\end{equation}
where $N_p$ is the number of particles, $\bm{x}_k^{(i)}$ is the state of particle $i$, $w_k^{(i)}$ is its normalized weight, and $\delta(\cdot)$ denotes the Dirac delta distribution. The weights satisfy
\begin{equation}
w_k^{(i)}\ge 0,
\qquad
\sum_{i=1}^{N_p}w_k^{(i)}=1.
\end{equation}

\subsection{Known-pose initialization}
For the current debugging baseline, the initial state is approximately known. The initial particle cloud is sampled independently as
\begin{equation}
\bm{x}_0^{(i)}\sim
\mathcal{N}\!\left(\bm{x}_0,\operatorname{diag}(\bm{\sigma}_0^2)\right),
\end{equation}
where $\bm{x}_0$ is the nominal initial state and $\bm{\sigma}_0$ is the vector of component-wise initial standard deviations. All particles initially have equal weight $w_0^{(i)}=1/N_p$ before the first range update.

\subsection{Particle propagation}
Particle $i$ is propagated from the measured IMU signal with additional process perturbations,
\begin{align}
\bm{a}_{k}^{b,(i)} &= \bm{a}_{m,k}^{b}+\bm{\epsilon}_{a,k}^{(i)},\\
\omega_{k}^{(i)} &= \omega_{m,k}+\epsilon_{\omega,k}^{(i)},
\end{align}
where
\begin{align}
\bm{\epsilon}_{a,k}^{(i)} &\sim \mathcal{N}(\bm{0},\sigma_{a,\mathrm{PF}}^2I_2),\\
\epsilon_{\omega,k}^{(i)} &\sim \mathcal{N}(0,\sigma_{\omega,\mathrm{PF}}^2).
\end{align}
The quantities $\sigma_{a,\mathrm{PF}}$ and $\sigma_{\omega,\mathrm{PF}}$ are design parameters of the PF process model. Each particle is then propagated with the same midpoint-yaw mechanization in Eqs.~\eqref{eq:mechanization_p}--\eqref{eq:mechanization_psi}, replacing the measured acceleration and yaw rate by the particle-specific values above.

\subsection{Range likelihood and weight update}
The position component of particle $i$ is denoted by $\bm{p}_k^{(i)}$. Its predicted range to the auxiliary node is
\begin{equation}
\hat d_k^{(i)}=\lVert\bm{p}_k^{(i)}-\bm{p}_{A,k}\rVert_2,
\end{equation}
where $\hat d_k^{(i)}$ is the particle-predicted range. The corresponding residual is
\begin{equation}
r_k^{(i)}=z_k-\hat d_k^{(i)}.
\end{equation}
Under the Gaussian UWB model in Eq.~\eqref{eq:uwb_measurement}, the likelihood is
\begin{equation}
p(z_k\mid\bm{x}_k^{(i)})
=\frac{1}{\sqrt{2\pi}\sigma_{\mathrm{UWB}}}
\exp\!\left[-\frac{(r_k^{(i)})^2}{2\sigma_{\mathrm{UWB}}^2}\right].
\end{equation}
The unnormalized importance weight is therefore
\begin{equation}
\tilde w_k^{(i)}=w_{k-1}^{(i)}p(z_k\mid\bm{x}_k^{(i)}),
\end{equation}
followed by normalization
\begin{equation}
w_k^{(i)}=\frac{\tilde w_k^{(i)}}{\sum_{j=1}^{N_p}\tilde w_k^{(j)}}.
\end{equation}
The implementation performs the likelihood calculation in log space for numerical robustness.

\subsection{State estimate and resampling}
The estimated position and velocity are the weighted particle means,
\begin{equation}
\hat{\bm{p}}_k=\sum_{i=1}^{N_p}w_k^{(i)}\bm{p}_k^{(i)},
\qquad
\hat{\bm{v}}_k=\sum_{i=1}^{N_p}w_k^{(i)}\bm{v}_k^{(i)}.
\end{equation}
Because yaw is an angle, it is estimated with a circular mean,
\begin{equation}
\hat\psi_k=\operatorname{atan2}\!\left(
\sum_{i=1}^{N_p}w_k^{(i)}\sin\psi_k^{(i)},
\sum_{i=1}^{N_p}w_k^{(i)}\cos\psi_k^{(i)}
\right).
\end{equation}
Together these quantities form the PF state estimate $\hat{\bm{x}}_k$.

Weight degeneracy is monitored using the effective sample size
\begin{equation}
N_{\mathrm{eff},k}=\frac{1}{\sum_{i=1}^{N_p}(w_k^{(i)})^2}.
\end{equation}
If
\begin{equation}
N_{\mathrm{eff},k}<\gamma N_p,
\end{equation}
where $\gamma\in(0,1]$ is the resampling fraction, systematic resampling is applied and the new particle weights are reset to $1/N_p$.

\section{Evaluation metrics}
Let $N_s$ denote the number of time samples included in an evaluation. The Euclidean position error at sample $k$ is
\begin{equation}
e_{p,k}=\lVert\hat{\bm{p}}_k-\bm{p}_k\rVert_2.
\end{equation}
The position root-mean-square error (RMSE) is
\begin{equation}
\mathrm{RMSE}_p=
\sqrt{\frac{1}{N_s}\sum_{k=0}^{N_s-1}e_{p,k}^2}.
\end{equation}
The 95th-percentile position error, abbreviated as p95 in result tables, is
\begin{equation}
e_{p,95}=Q_{0.95}\!\left(\{e_{p,k}\}_{k=0}^{N_s-1}\right),
\end{equation}
where $Q_{0.95}$ is the empirical 95th-quantile operator. The final position error is simply $e_{p,N_s-1}$.

The wrapped yaw error is
\begin{equation}
e_{\psi,k}=\operatorname{wrap}(\hat\psi_k-\psi_k),
\end{equation}
and the yaw RMSE is
\begin{equation}
\mathrm{RMSE}_{\psi}=
\sqrt{\frac{1}{N_s}\sum_{k=0}^{N_s-1}e_{\psi,k}^2}.
\end{equation}
Yaw RMSE is converted from radians to degrees when reported in the current result tables.

\section{Local observability diagnostic}
The current observability check uses the ideal, bias-free five-state model from Eq.~\eqref{eq:state}. Define
\begin{equation}
\bm{a}_k^n=R_b^n(\psi_{k+1/2})\bm{a}_k^b
\end{equation}
and its yaw sensitivity
\begin{equation}
\bm{a}_{\psi,k}=\frac{\partial\bm{a}_k^n}{\partial\psi}.
\end{equation}
The local state-transition Jacobian $F_k=\partial\bm{x}_{k+1}/\partial\bm{x}_k$ used in the current diagnostic is
\begin{equation}
F_k=
\begin{bmatrix}
I_2 & \Delta t I_2 & \frac12\Delta t^2\bm{a}_{\psi,k}\\
0_{2\times2} & I_2 & \Delta t\bm{a}_{\psi,k}\\
0_{1\times2} & 0_{1\times2} & 1
\end{bmatrix},
\end{equation}
where $I_2$ is the $2\times2$ identity matrix and $0_{r\times c}$ denotes an $r\times c$ zero matrix.

For the scalar range measurement, the measurement Jacobian $H_k=\partial d_k/\partial\bm{x}_k$ is
\begin{equation}
H_k=
\begin{bmatrix}
\dfrac{(\bm{p}_k-\bm{p}_{A,k})^\top}{d_k} & 0 & 0 & 0
\end{bmatrix}.
\end{equation}
For $j>0$, define the state-transition sensitivity from sample $0$ to $j$ as
\begin{equation}
\Phi_{j,0}=F_{j-1}F_{j-2}\cdots F_0,
\qquad
\Phi_{0,0}=I_5.
\end{equation}
The finite-horizon local observability matrix from sample $0$ through $K$ is then
\begin{equation}
\mathcal{O}_{0:K}=
\begin{bmatrix}
H_0\\
H_1\Phi_{1,0}\\
H_2\Phi_{2,0}\\
\vdots\\
H_K\Phi_{K,0}
\end{bmatrix}.
\end{equation}
Its singular values are inspected numerically. A zero smallest singular value $\sigma_{\min}(\mathcal{O}_{0:K})$ indicates a locally unobservable direction in this linearized finite-horizon diagnostic. The ratio
\begin{equation}
\frac{\sigma_{\min}(\mathcal{O}_{0:K})}{\sigma_{\max}(\mathcal{O}_{0:K})}
\end{equation}
is used only as a conditioning indicator. This test is not presented as a complete nonlinear observability proof.

\section{Relation to the $\Delta L$, $\Delta\psi$ model in Han et al.}
Han et al.~\cite{han2020cooperative} formulate their cooperative layer with the compact state $[p_x,p_y,\psi]^\top$. Their filter receives a travelled-distance increment $\Delta L$ and an azimuth/yaw increment $\Delta\psi$ from an underlying DR/INS system. The initial repository bootstrap used that high-level interface directly. The revised thesis implementation moves one level closer to the physical sensor by explicitly generating and integrating acceleration and yaw-rate measurements, which requires velocity to be part of the state. The later AACOPF reproduction must therefore distinguish clearly between (i) the original $\Delta L,\Delta\psi$ interface used for a faithful paper baseline and (ii) the IMU-driven five-state formulation used for the hardware-oriented thesis development.
